\documentclass[10pt,letterpaper]{article}

\usepackage[T1]{fontenc}
\usepackage[utf8]{inputenc}
\usepackage[top=0.62in,bottom=0.62in,left=0.7in,right=0.7in]{geometry}
\usepackage{titlesec}
\usepackage{enumitem}
\usepackage[hidelinks]{hyperref}

\pagestyle{empty}
\setlength{\parindent}{0pt}

% Section formatting matching Yixuan Even Xu's Computer Modern Small-Caps CV style
\titleformat{\section}
  {\large\scshape\raggedright}
  {}{0em}
  {}
  [\titlerule]
\titlespacing*{\section}{0pt}{9pt}{5pt}

% Compact bullet list matching 10pt hanging indent
\setlist[itemize]{
  leftmargin=10pt,
  label=\textbullet,
  itemsep=3.5pt,
  parsep=1pt,
  topsep=1.5pt
}

\begin{document}

% ==================== HEADER ====================
\begin{center}
  {\Huge\scshape Huang Yuchang} \\[6pt]
  \texttt{huangyuchang@tju.edu.cn}
  \textbf{~$\cdot$~}
  \texttt{huang-zirun.github.io}
  \textbf{~$\cdot$~}
  \texttt{github.com/huang-zirun} \\[2pt]
  International Joint Institute, Fuzhou \textbf{~$\cdot$~} Tianjin University \textbf{~$\cdot$~} China \textbf{~$\cdot$~} Last Update: November 16, 2026
\end{center}

\vspace{-2pt}

% ==================== EDUCATION ====================
\section{Education}

\textbf{Tianjin University, Fuzhou / Tianjin, China} \hfill 2025.09 -- Present \\
M.Eng. in Electronic Information (International Joint Institute, Fuzhou) \\
Research \& Engineering Focus: AI for Software Engineering (AI4SE), AI for Science (AI4Science)

% ==================== PROFESSIONAL EXPERIENCE ====================
\section{Professional Experience}

\begin{itemize}
  \item \textbf{ByteDance Inc. (Feishu / Lark)} \hfill 2026.06 -- Present \\
  R\&D Intern. Building and maintaining the team's end-to-end requirement delivery harness, and designing, developing, and iterating on internal AI agents to streamline engineering workflows and automated code delivery.
\end{itemize}

% ==================== PROJECTS ====================
\section{Selected Projects}

\begin{itemize}
  \item \textbf{DearAlpha: Intelligent Quantitative Alpha Factor Mining Framework} \hfill 2026.04 -- 2026.09 \\
  \emph{Tech Stack: Python, WorldQuant Brain API, Optuna TPE, Bayesian Optimization, Multi-LLM, Pandas} \\
  Built around an ``AI-driven hypothesis generation, programmatic verification, and human final judgment'' architecture (\texttt{huang-zirun/DearAlpha}). Uses multi-backend LLMs guided by economic themes (momentum, value, volatility) to synthesize template-free raw alpha expressions; integrates a 3-stage progressive expansion pipeline, two-pass layered pruning (cutting redundant backtest simulations by $\sim$50\%), Optuna TPE Bayesian hyperparameter optimization, and multi-metric quality gates (Sharpe, Fitness, Turnover, and correlation filtering).

  \item \textbf{DevFlow Engine: AI-Driven End-to-End Requirement-to-Delivery Pipeline} \hfill 2026.04 -- 2026.05 \\
  \emph{Tech Stack: Python, LangGraph, Python AST, Tree-sitter, lark-cli WebSocket, JSON Schema, OpenAPI 3.0} \\
  Designed a ``deterministic pipeline skeleton, specialized agent execution, and human-in-the-loop governance'' architecture (\texttt{huang-zirun/lark-agent}) that decomposes the software delivery lifecycle into a 6-stage state graph (intake, solution design, code generation, testing, code review, and delivery packaging) decoupled via versioned JSON contracts. Combines incremental syntax-tree semantic indexing with workspace-sandboxed toolchains, a bounded review-to-code self-repair loop, and dual Feishu approval checkpoints.
\end{itemize}

% ==================== COMPETITIONS AND AWARDS ====================
\section{Honors and Awards}

\begin{itemize}
  \item \textbf{2026 Global AI Innovation, Governance, and Safety Competition, Second Prize} \hfill 2026.11 \\
  AI Agent Application Track (Project: \emph{Maestro AI}; \textbf{10 track finalists} / \textbf{30 finalists} out of 118 teams worldwide), RUC

  \item \textbf{The 2nd Alibaba ``AI for Good'' Innovation Challenge, Open Source Contribution Award} \hfill 2026.09 \\
  Rural Education Track (Project: \emph{Sonara}; \textbf{10 teams nationwide} out of 1,111 registered teams), Alibaba Group

  \item \textbf{2026 Feishu AI Campus Challenge, AI Rising Star Certificate \& Final Interview Fast-Track} \hfill 2026.05 \\
  Track: AI-Driven Requirement-to-Delivery Pipeline Engine, hosted by ByteDance (Feishu / Lark)

  \item \textbf{The 17th Chinese Mathematics Competition (CMC), First Prize in Guangdong Province} \hfill 2025.11 \\
  Awarded to Top 5\% participants in Guangdong Province, hosted by Chinese Mathematical Society (CMS)

  \item \textbf{Contemporary Undergraduate Mathematical Contest in Modeling (CUMCM), Second Prize} \hfill 2023.09 \\
  Guangdong Province, hosted by China Society for Industrial and Applied Mathematics (CSIAM) (Team Member)

  \item \textbf{GE Lean \& Green Challenge, National Champion} \hfill 2023.07 \\
  Awarded 1st place nationwide, hosted by General Electric (GE) (Team Member)

  \item \textbf{Contemporary Undergraduate Mathematical Contest in Modeling (CUMCM), Second Prize} \hfill 2022.09 \\
  Guangdong Province, hosted by China Society for Industrial and Applied Mathematics (CSIAM) (Team Member)
\end{itemize}

% ==================== TECHNICAL SKILLS ====================
\section{Technical Skills}

\begin{itemize}
  \item \textbf{Systems \& Agent Engineering}: Python, LangGraph, Multi-Agent Harness Design, AST / Tree-sitter Codebase Indexing, JSON Schema Contract Design, OpenAPI 3.0, Feishu / Lark Open Platform (\texttt{lark-cli})
  \item \textbf{AI for Science \& Quantitative Modeling}: Multi-LLM Orchestration, Optuna TPE Bayesian Optimization, Quantitative Alpha Factor Mining, Pandas, Mathematical Modeling
\end{itemize}

\end{document}
